\documentclass[11pt]{article}
\usepackage[margin=1in]{geometry}
\usepackage{amsmath,amssymb,amsthm}
\usepackage{graphicx}
\usepackage{hyperref}

\newtheorem{theorem}{Theorem}
\newtheorem{proposition}{Proposition}
\newtheorem{lemma}{Lemma}

\newcommand{\E}{\mathbb E}
\newcommand{\Pp}{\mathbb P}
\newcommand{\norm}[1]{\left\lVert #1\right\rVert}

\title{Hilbert-Valued Characteristic Domination for \(p\geq 2\)}
\author{Agent \texttt{agent\_lane\_19}}
\date{2026-06-16}

\begin{document}
\maketitle

\section*{Status}

This packet gives a candidate partial result, likely valid, for Remark 12.10 of
Yaroslavtsev's paper \emph{Local characteristics and tangency of vector-valued
martingales} \cite{Yaroslavtsev2019}.  It strengthens the earlier Hilbert
\(p=2\) packet to every exponent \(p\geq2\).  The result is still partial
relative to the source problem, because it does not cover non-Hilbert UMD
spaces or \(1\leq p<2\).

\section*{Source Problem}

The source is arXiv:1907.11588, PDF page 105, Remark 12.10.  The source crop
below states the open general-local-martingale issue and its discrete
independent-increment reduction.

\begin{center}
\includegraphics[width=\textwidth]{figures/open_problem_crop.png}
\end{center}

The next page adds that, by symmetrization, one can assume the variables are
symmetric, and that even this case was not known to the author.

\begin{center}

\includegraphics[width=\textwidth]{figures/open_problem_continuation_crop.png}
\end{center}

In discrete form, the source asks whether independent mean-zero \(X\)-valued
variables \((\xi_n)\), \((\eta_n)\) satisfying
\[
   \sum_n \Pp(\xi_n\in B)
   \leq
   \sum_n \Pp(\eta_n\in B)
   \qquad
   (B\subset X\setminus\{0\}\text{ Borel})
\]
must satisfy
\[
   \E\norm{\sum_n \xi_n}^p
   \leq
   C_{p,X}\E\norm{\sum_n \eta_n}^p .
\]

\section*{Idea of the Proof}

The \(c_0\) counterexample works because the norm sees the largest occupancy
after the same aggregate jump mass is regrouped.  Hilbert spaces do not have
that freedom at exponents \(p\geq2\).  A Hilbert-valued Rosenthal inequality
says that the \(p\)-th moment of an independent mean-zero sum is equivalent to
only two aggregate quantities:
\[
  \left(\sum_n \E\norm{Z_n}^2\right)^{p/2}
  +
  \sum_n \E\norm{Z_n}^p .
\]
Both are monotone under the aggregate measure domination, by integrating
\(\norm{x}^2\) and \(\norm{x}^p\).  This proves the discrete comparison in
Hilbert space for \(p\geq2\).  The source paper already reduces the missing
accessible-jump part of the general local martingale problem to precisely this
discrete comparison; the continuous and quasi-left-continuous parts are already
handled there.

\section*{Discrete Hilbert Comparison}

\begin{lemma}[Hilbert Rosenthal estimate]\label{lem:hilbert-rosenthal}
Let \(H\) be a Hilbert space, \(p\geq2\), and let
\(Z_1,\ldots,Z_N\) be independent mean-zero \(H\)-valued random variables.
Then
\[
 \E\norm{\sum_{n=1}^N Z_n}^p
 \eqsim_p
 \left(\sum_{n=1}^N \E\norm{Z_n}^2\right)^{p/2}
 +
 \sum_{n=1}^N \E\norm{Z_n}^p .
\]
\end{lemma}

\begin{proof}
Let \(S_k=\sum_{n=1}^k Z_n\).  The Hilbert-valued discrete
Burkholder-Davis-Gundy inequality gives
\[
  \E\sup_{1\leq k\leq N}\norm{S_k}^p
  \eqsim_p
  \E\left(\sum_{n=1}^N\norm{Z_n}^2\right)^{p/2};
\]
this is the Hilbert case of the discrete estimate recalled in
\cite{Yaroslavtsev2018}.  Since \((\norm{S_k})_k\) is a submartingale, Doob's
\(L^p\) inequality also gives
\[
  \E\sup_{1\leq k\leq N}\norm{S_k}^p
  \lesssim_p
  \E\norm{S_N}^p,
\]
while the reverse inequality is immediate.  Therefore
\[
  \E\norm{S_N}^p
  \eqsim_p
  \E\left(\sum_{n=1}^N\norm{Z_n}^2\right)^{p/2}.
\]
Finally apply the scalar Rosenthal inequality to the independent nonnegative
variables \(Y_n=\norm{Z_n}^2\) with exponent \(p/2\geq1\):
\[
 \E\left(\sum_{n=1}^N Y_n\right)^{p/2}
 \eqsim_p
 \left(\sum_{n=1}^N \E Y_n\right)^{p/2}
 +
 \sum_{n=1}^N \E Y_n^{p/2}.
\]
This is the displayed estimate.
\end{proof}

\begin{proposition}[Aggregate domination in Hilbert spaces for \(p\geq2\)]
\label{prop:discrete-hilbert}
Let \(H\) be a Hilbert space and \(p\geq2\).  Let
\((\xi_n)_{n=1}^N\) and \((\eta_n)_{n=1}^N\) be independent mean-zero
\(H\)-valued random variables.  If
\[
   \sum_{n=1}^N \Pp(\xi_n\in B)
   \leq
   \sum_{n=1}^N \Pp(\eta_n\in B)
   \qquad
   (B\subset H\setminus\{0\}\text{ Borel}),
\]
then
\[
   \E\norm{\sum_{n=1}^N \xi_n}^p
   \leq
   C_p\,\E\norm{\sum_{n=1}^N \eta_n}^p .
\]
\end{proposition}

\begin{proof}
For \(q=2\) and \(q=p\), integrate \(\norm{x}^q\) against the two aggregate
measures.  The atom at \(0\) is irrelevant.  We obtain
\[
  \sum_{n=1}^N\E\norm{\xi_n}^q
  \leq
  \sum_{n=1}^N\E\norm{\eta_n}^q,
  \qquad q\in\{2,p\}.
\]
Applying Lemma \ref{lem:hilbert-rosenthal} to \((\xi_n)\), using these two
inequalities, and applying Lemma \ref{lem:hilbert-rosenthal} again to
\((\eta_n)\), gives the result.
\end{proof}

\section*{General Local Martingales}

\begin{theorem}
Let \(H\) be a separable Hilbert space and let \(p\geq2\).  Let \(M\) and \(N\)
be \(H\)-valued local martingales such that \(N\) is characteristically
dominated by \(M\) in the sense of \cite{Yaroslavtsev2019}.  Then
\[
  \E\sup_{t\geq0}\norm{N_t}^p
  \leq
  C_p\,\E\sup_{t\geq0}\norm{M_t}^p ,
\]
whenever the right hand side is finite.
\end{theorem}

\begin{proof}
We use the reduction developed in the source paper.  The continuous part is
handled by the characteristic-subordination/continuous domination estimate
used in the proof of Theorem 12.8 of \cite{Yaroslavtsev2019}.  The
quasi-left-continuous purely discontinuous part is exactly Theorem 12.8 of
\cite{Yaroslavtsev2019}.  Thus, after the canonical decomposition, the only
missing piece is the accessible-jump part.

For accessible jumps, use the accessible-jump decoupled tangent construction
and the accessible-jump approximation proposition cited explicitly in Remark
12.10 of \cite{Yaroslavtsev2019}.  These approximate by martingales with
finitely many predictable jumps and represent the corresponding decoupled
tangent object, for almost every base point, as a sum of independent mean-zero
increments.  In Remark 12.10 the source states that the main remaining issue
for general characteristic domination is precisely the discrete
independent-increment aggregate-domination problem.  Under characteristic
domination, the total compensator of the accessible jumps of \(N\) is dominated
by that of \(M\); for the decoupled tangent sums this is exactly the aggregate
domination hypothesis in Proposition \ref{prop:discrete-hilbert}.  Applying
that proposition pointwise to the decoupled accessible-jump sums gives the
desired \(L^p\) comparison for the accessible parts.  The standard
tangent/decoupling estimate for Hilbert spaces, recalled as Theorem 1.1 in
\cite{Yaroslavtsev2019}, transfers the estimate back to the original
accessible-jump martingales, and the approximation then removes the
finite-jump truncation.

Combining the continuous, quasi-left-continuous, and accessible estimates with
the canonical-decomposition \(L^p\) bounds used in the source proof yields the
displayed estimate.
\end{proof}

\section*{Verification Notes}

The new proof step is Proposition \ref{prop:discrete-hilbert}.  It is stable
under symmetrization but does not need symmetry: independence and the mean-zero
assumption are enough for the Hilbert Rosenthal estimate.  The proof also
explains why the \(c_0\) counterexample cannot be transplanted into Hilbert
space for \(p\geq2\): after Rosenthal, regrouping the same aggregate mass can
change neither the total second moment nor the total \(p\)-th moment enough to
create an unbounded ratio.

\section*{Novelty and Literature Check}

The cheap run indexes were searched for the arXiv id and for the terms
\texttt{characteristically dominated}, \texttt{aggregate domination},
\texttt{Hilbert}, \texttt{p>=2}, \texttt{Rosenthal}, and
\texttt{Bichteler-Jacod}.  The run already contained the broad \(c_0\)
counterexample and the Hilbert \(p=2\) packet, but no Hilbert \(p\geq2\)
packet.  Literature searches found Yaroslavtsev's BDG paper
\cite{Yaroslavtsev2018}, Rosenthal/Pinelis moment tools
\cite{Pinelis2012}, and Montgomery-Smith--Pruss comparison inequalities
\cite{MontgomerySmithPruss2001}, but no exact recorded solution of the Hilbert
\(p\geq2\) subcase of Remark 12.10.

\section*{Limitations and Next Routes}

This packet does not solve \(1\leq p<2\).  The Rosenthal reduction used above
is naturally aligned with \(p\geq2\), because the \(p\)-th moment controls both
the quadratic characteristic and the large-jump characteristic.  It also does
not solve general non-Hilbert UMD spaces.  A plausible next step is to test
whether the same two-characteristic reduction holds in \(2\)-smooth UMD spaces
or in \(L^q\)-spaces with the Dirksen--Yaroslavtsev/Marinelli--Roeckner
Burkholder-Novikov-Rosenthal estimates.

\section*{Human Review Recommendation}

First verify Lemma \ref{lem:hilbert-rosenthal}, especially the use of scalar
Rosenthal for \(\sum\norm{Z_n}^2\).  Then verify that the source paper's
accessible-jump reduction indeed needs only the discrete aggregate-domination
comparison; this is the point explicitly isolated in Remark 12.10.

\begin{thebibliography}{9}

\bibitem{Yaroslavtsev2019}
I. S. Yaroslavtsev,
\emph{Local characteristics and tangency of vector-valued martingales},
arXiv:1907.11588, 2019.

\bibitem{Yaroslavtsev2018}
I. S. Yaroslavtsev,
\emph{Burkholder-Davis-Gundy inequalities in UMD Banach spaces},
arXiv:1807.05573, 2018.

\bibitem{Pinelis2012}
I. Pinelis,
\emph{Rosenthal-type inequalities for martingales in 2-smooth Banach spaces},
arXiv:1212.1912, 2012.

\bibitem{MontgomerySmithPruss2001}
S. J. Montgomery-Smith and A. R. Pruss,
\emph{A comparison inequality for sums of independent random variables},
J. Math. Anal. Appl. 254 (2001), no. 1, 35--42;
arXiv:math/9811124.

\bibitem{Rosenthal1970}
H. P. Rosenthal,
\emph{On the subspaces of \(L^p\) \((p>2)\) spanned by sequences of
independent random variables},
Israel J. Math. 8 (1970), 273--303.

\end{thebibliography}

\end{document}
